\begin{exercise}{56}
We call a set $S$ \emph{semi-saturated} if it is closed under pushouts, retracts and
countable compositions. Show that a semi-saturated set
\begin{enumerate}[label=(\arabic*)]
\item contains isomorphisms and is closed under finite coproducts, if it contains the identity of
an initial object $\varnothing$;
\item is closed under composition, i.e., if $f\colon A\to B$ and $g\colon B\to C$ are elements of $S$, then
so is $gf\colon A\to C$;
\item is closed under countable coproducts if it is closed under finite coproducts, i.e., if
$\{f_i\colon A_i\to B_i\}_{i\in I}$ is a countable family of elements of $S$, then the map $\coprod_{i\in I}f_i\colon\allowbreak\coprod_{i\in I}A_i\to\coprod_{i\in I}B_i$\footnote{The original prints $\coprod_{i\in I}\colon A_i\to B_i$, omitting the coproduct operands. Restore $f_i$, $A_i$, and $B_i$ as displayed: the induced map has source $\coprod_iA_i$ and target $\coprod_iB_i$, and its restriction to $A_i$ is $f_i$ followed by the inclusion of $B_i$.} is
an element of $S$ as well.
\end{enumerate}

\begin{proof}
Work in $\mathrm{sSet}$ as in Definition~1.3.5, or in any
category admitting the colimits below.

(2) First form $P=C\amalg_B C$, the pushout of $g$ along itself.
One inclusion $j\colon C\to P$ belongs to $S$.
The fold $r\colon P\to C$ satisfies $rj=\mathrm{id}_C$,
so $\mathrm{id}_C$ is a retract of $j$ as an arrow
(use identity maps on their common domain $C$).
Thus $\mathrm{id}_C\in S$ without the hypothesis of (1).
The sequence
\[
 A\xrightarrow{f}B\xrightarrow{g}C
 \xrightarrow{\mathrm{id}_C}C\xrightarrow{\mathrm{id}_C}\cdots
\]
has colimit $C$ and initial structure map $gf$.
Closure under countable compositions proves (2).

(1) Every isomorphism is a pushout of
$\mathrm{id}_{\varnothing}$ along the map from $\varnothing$
to its domain, so it lies in $S$.
A finite coproduct of maps in $S$ factors into finitely many
maps replacing one source summand by its target.
Each replacement is a pushout of that summand's map,
so (2) puts their composite in $S$.
The empty coproduct is $\mathrm{id}_{\varnothing}$.

(3) Finite families are covered by the hypothesis.
For an infinite countable family, enumerate it by $i\in\mathbb N$
and put
\[
 Z_k=\Bigl(\coprod_{i<k}B_i\Bigr)
       \amalg\Bigl(\coprod_{i\ge k}A_i\Bigr).
\]
The transition $Z_k\to Z_{k+1}$ replaces $A_k$ by $B_k$
using $f_k$ and fixes all other summands.
It is a pushout of $f_k$, hence belongs to $S$.
A cocone from this sequence to $T$ is uniquely determined
by its maps $B_i\to T$ once the $i$th summand has been replaced;
earlier values on $A_i$ are then their composites with $f_i$.
Consequently its colimit is $\coprod_iB_i$,
and the map from $Z_0=\coprod_iA_i$ is $\coprod_i f_i$.
Countable-composition closure finishes the proof.
\end{proof}
\end{exercise}
